\documentclass[10pt]{article}
\usepackage[margin=0.72in]{geometry}
\usepackage{amsmath,amssymb}
\usepackage{booktabs}
\usepackage{array}
\usepackage{enumitem}
\usepackage{hyperref}
\usepackage{xcolor}
\usepackage{titlesec}
\usepackage{listings}
\usepackage{microtype}

\titleformat{\section}{\Large\bfseries}{\thesection}{0.7em}{}
\titleformat{\subsection}{\large\bfseries}{\thesubsection}{0.7em}{}

\lstset{
  basicstyle=\ttfamily\footnotesize,
  columns=fullflexible,
  keepspaces=true,
  frame=single,
  breaklines=true,
  showstringspaces=false
}

\newcommand{\bigO}{O}
\newcommand{\SA}{\operatorname{SA}}
\newcommand{\LCP}{\operatorname{LCP}}
\newcommand{\rankv}{\operatorname{rank}}
\newcommand{\lca}{\operatorname{lca}}
\newcommand{\dist}{\operatorname{dist}}
\newcommand{\RMQ}{\operatorname{RMQ}}
\newcommand{\pref}{\operatorname{pref}}
\newcommand{\zfunc}{\operatorname{Z}}

\title{Problemas Racso de Static Trees y Strings\\
\large Enunciados interpretados, teoria aplicada, soluciones de examen y pseudocodigo}
\author{}
\date{}

\begin{document}
\sloppy
\maketitle

\section*{Como leer este PDF}
Este documento toma los problemas de \texttt{soluciones\_teoria\_racso.pdf} que
corresponden a la carpeta \texttt{statics and strings/}: RMQ, LCA/static trees y
strings. No copia literalmente los enunciados de los jueces: cada problema tiene
un \textbf{enunciado interpretado en espanol}, suficiente para estudiar y
responder en papel sin depender de memorizar texto.

El objetivo es que aprendas teoria mediante problemas. En cada seccion veras:
\begin{itemize}[leftmargin=*]
\item como interpretar el enunciado;
\item que teoria debes activar;
\item que estructura usar y por que;
\item algoritmo paso a paso;
\item pseudocodigo corto estilo C++;
\item correctitud;
\item complejidad;
\item trampas de examen.
\end{itemize}

\subsection*{Problemas incluidos}
\begin{center}
\begin{tabular}{p{0.33\linewidth}p{0.45\linewidth}p{0.14\linewidth}}
\toprule
Tema & Problemas & Prioridad \\
\midrule
RMQ/static array & Static RMQ & Maxima \\
Strings base & Suffix array y LCP, busqueda, conteo & Maxima \\
Strings clasicos & distintas, LCS, ordenar substrings, bordes & Maxima/alta \\
Strings avanzados & Refran, substring periodica, Frequency of String & Alta \\
Static trees/LCA & Imperial Roads, Fools and Roads & Maxima/alta \\
\bottomrule
\end{tabular}
\end{center}

\tableofcontents
\newpage

\section{Base teorica comun: lo minimo que debes tener en la cabeza}

\subsection{Static significa que puedes preprocesar}
Si el arreglo, arbol o string no cambia, puedes gastar tiempo antes de las
consultas. Eso se llama preprocesamiento. En problemas de examen, la frase
``muchas consultas'' casi siempre significa:
\begin{quote}
No resuelvas cada consulta desde cero; construye una estructura.
\end{quote}

Ejemplos:
\begin{itemize}[leftmargin=*]
\item RMQ estatico: sparse table.
\item LCA en arbol fijo: binary lifting o Euler tour + RMQ.
\item Texto fijo con muchos patrones: suffix array o Aho-Corasick.
\item Comparar substrings muchas veces: suffix array + LCP + RMQ.
\end{itemize}

\subsection{RMQ}
RMQ significa Range Minimum Query. Dado un arreglo fijo $A$, una consulta
$\RMQ(l,r)$ pide el minimo valor o la posicion del minimo en $A[l..r]$.

La estructura mas directa es sparse table:
\[
st[k][i]=\min(A[i],A[i+1],\ldots,A[i+2^k-1]).
\]
Para consultar $[l,r]$, tomas $k=\lfloor\log_2(r-l+1)\rfloor$ y combinas:
\[
\min(st[k][l],st[k][r-2^k+1]).
\]
Funciona porque los dos bloques cubren el rango y el minimo es idempotente:
$\min(x,x)=x$.

\subsection{LCA}
El Lowest Common Ancestor de $u$ y $v$ es el ancestro comun mas profundo. Sirve
para rutas en arbol:
\[
\dist(u,v)=depth[u]+depth[v]-2depth[\lca(u,v)].
\]
Si ademas guardas maximos/minimos en saltos de binary lifting, puedes responder
max edge on path o min edge on path.

\subsection{Suffix array, rank y LCP}
El suffix array de $S$ ordena todos los sufijos:
\[
S[\SA[0]..] < S[\SA[1]..] < \cdots < S[\SA[n-1]..].
\]
El rank inverso dice en que posicion aparece cada sufijo:
\[
\SA[\rankv[i]]=i.
\]
El LCP array guarda:
\[
\LCP[i]=lcp(S[\SA[i-1]..],S[\SA[i]..]).
\]
Si dos sufijos estan en posiciones $a<b$ del suffix array:
\[
lcp(i,j)=\min(\LCP[a+1],\ldots,\LCP[b]).
\]
Por eso se preprocesa RMQ sobre LCP para comparar substrings en $\bigO(1)$.

\subsection{Como clasificar rapido}
\begin{center}
\begin{tabular}{p{0.37\linewidth}p{0.42\linewidth}}
\toprule
Frase del problema & Estructura que debes pensar \\
\midrule
minimo en rango, sin updates & sparse table / RMQ \\
camino en arbol, distancia, max edge & LCA / binary lifting augmentado \\
muchos caminos y contar aristas usadas & LCA + diferencia en arbol \\
buscar patron en texto fijo & suffix array, intervalo de sufijos \\
subcadenas distintas & total menos suma de LCP \\
subcadena comun mas larga & SA de $A\#B$ + LCP \\
comparar/ordenar substrings & rank + LCP arbitrario + RMQ \\
borde, prefijo=sufijo & KMP prefix function o Z-function \\
periodicidad & KMP/Z o igualdad por LCP \\
muchos patrones en un texto & Aho-Corasick o suffix array segun enfoque \\
\bottomrule
\end{tabular}
\end{center}

\newpage
\section{Static RMQ}
\textbf{Link.} \url{https://codeforces.com/group/AmklRC8lhZ/contest/694840/problem/A}

\subsection{Enunciado interpretado}
Se te da un arreglo fijo de $n$ enteros. Luego llegan muchas consultas
$(l,r)$. Para cada una debes responder el minimo del subarreglo desde $l$ hasta
$r$. El arreglo no recibe updates.

\subsection{Como interpretarlo en examen}
La palabra clave es \textbf{static}. Eso descarta segment tree como primera
opcion si solo hay minimos y no hay updates. Lo correcto es decir: preproceso
bloques de potencias de dos con sparse table y respondo cada rango usando dos
bloques solapados.

\subsection{Teoria necesaria}
Sparse table guarda respuestas para rangos canonicos. Canonico significa: rango
que empieza en alguna posicion $i$ y tiene longitud $2^k$. La razon de elegir
potencias de dos es que cualquier longitud contiene una potencia de dos maxima.

Para $[l,r]$, si $2^k\le r-l+1$, entonces:
\[
[l,l+2^k-1]\quad\text{y}\quad[r-2^k+1,r]
\]
cubren todo el intervalo. Se pueden solapar. Para minimo eso no afecta, porque
repetir elementos no cambia la respuesta.

\subsection{Algoritmo}
\begin{enumerate}[leftmargin=*]
\item Precalcula $lg[len]=\lfloor\log_2 len\rfloor$.
\item Inicializa $st[0][i]=A[i]$.
\item Para cada $k>0$, calcula:
\[
st[k][i]=\min(st[k-1][i],st[k-1][i+2^{k-1}]).
\]
\item Para query $(l,r)$:
\[
k=lg[r-l+1],
\]
\[
ans=\min(st[k][l],st[k][r-2^k+1]).
\]
\end{enumerate}

\subsection{Pseudocodigo corto}
\begin{lstlisting}[language=C++]
build(A):
  lg[1] = 0
  for i = 2..n: lg[i] = lg[i/2] + 1
  for i = 0..n-1: st[0][i] = A[i]
  for k = 1..lg[n]:
    for i = 0; i + (1<<k) <= n; i++:
      st[k][i] = min(st[k-1][i],
                     st[k-1][i + (1<<(k-1))])

query(l, r):
  k = lg[r-l+1]
  return min(st[k][l], st[k][r-(1<<k)+1])
\end{lstlisting}

\subsection{Correctitud}
El preprocesamiento es correcto por induccion en $k$. Para $k=0$, cada bloque
tiene un solo elemento. Para $k>0$, un bloque de longitud $2^k$ es la union de
dos bloques de longitud $2^{k-1}$, por lo que tomar el minimo de sus respuestas
da el minimo del bloque grande.

En la consulta, los dos bloques elegidos cubren $[l,r]$. Todo elemento del rango
esta en al menos uno de ellos, asi que el minimo global esta considerado. Si un
elemento esta en ambos, no afecta por idempotencia del minimo. Por tanto la
respuesta es exactamente el minimo del rango.

\subsection{Complejidad y trampas}
Preprocesamiento $\bigO(n\log n)$, memoria $\bigO(n\log n)$, query $\bigO(1)$.
Trampa: sparse table sirve para min, max y gcd; no sirve asi para suma porque
dos bloques solapados contarian dos veces.

\newpage
\section{Suffix array y LCP array}
\textbf{Link.} \url{https://codeforces.com/group/AmklRC8lhZ/contest/695758/problem/A}

\subsection{Enunciado interpretado}
Se te da una cadena $S$. Debes construir o razonar sobre su suffix array y su
LCP array. Este problema es el fundamento de casi todos los demas de strings.

\subsection{Como interpretarlo en examen}
Si te piden suffix array, no pienses en comparar todos los sufijos con
comparacion lenta. Debes explicar una construccion: Manber-Myers si quieres una
respuesta simple, o DC3/SA-IS si piden linealidad. Para el examen, Manber-Myers
es la ruta mas defendible.

\subsection{Teoria necesaria}
Un sufijo es $S[i..n-1]$. El suffix array es una permutacion de posiciones. No
guarda strings; guarda indices. Si $S=\texttt{banana\$}$, el suffix array
ordenado es:
\[
\SA=[6,5,3,1,0,4,2].
\]
El LCP entre sufijos vecinos mide cuanto comparten al inicio. En
\texttt{banana\$}, \texttt{ana\$} y \texttt{anana\$} tienen LCP $3$.

\subsection{Algoritmo}
Manber-Myers ordena por prefijos de longitud $1,2,4,8,\ldots$.
\begin{enumerate}[leftmargin=*]
\item Ordena posiciones por caracter y asigna ranks iniciales.
\item En cada ronda de longitud $len$, representa el sufijo que empieza en $i$
por el par:
\[
(rank[i],rank[i+len]).
\]
\item Ordena esos pares.
\item Reasigna ranks: pares iguales reciben el mismo rank, pares distintos
ranks crecientes.
\item Duplica $len$ hasta cubrir toda la cadena.
\item Construye LCP con Kasai usando rank inverso.
\end{enumerate}

\subsection{Pseudocodigo corto}
\begin{lstlisting}[language=C++]
build_suffix_array(S):
  n = S.size()
  sa = [0,1,...,n-1]
  rank[i] = S[i]
  for len = 1; len < n; len *= 2:
    sort sa by (rank[i], rank[i+len] or -1)
    newRank[sa[0]] = 0
    for t = 1..n-1:
      newRank[sa[t]] = newRank[sa[t-1]]
      if pair(sa[t]) != pair(sa[t-1]):
        newRank[sa[t]]++
    rank = newRank
  return sa

build_lcp(S, sa):
  rank[sa[i]] = i
  h = 0
  for i = 0..n-1:
    if rank[i] == 0: continue
    j = sa[rank[i]-1]
    while S[i+h] == S[j+h]: h++
    lcp[rank[i]] = h
    if h > 0: h--
\end{lstlisting}

\subsection{Correctitud}
La induccion es la parte importante. Al inicio, $rank[i]$ ordena prefijos de
longitud $1$. Si en una ronda $rank$ ordena prefijos de longitud $len$, entonces
el par $(rank[i],rank[i+len])$ ordena prefijos de longitud $2len$: primero mira
la mitad izquierda y si empata mira la mitad derecha. Al duplicar hasta
$len\ge n$, ordenar prefijos equivale a ordenar sufijos completos.

Kasai es correcto porque al pasar de sufijo $i$ a sufijo $i+1$, el LCP con su
vecino lexicografico disminuye como mucho en uno; por eso reutiliza $h$ y suma
tiempo lineal total.

\subsection{Complejidad y trampas}
Manber-Myers con sort general cuesta $\bigO(n\log^2 n)$; con counting/radix sort
por ranks cuesta $\bigO(n\log n)$. Kasai cuesta $\bigO(n)$. Trampa de examen:
no confundas $\SA$ con $rank$; $\SA[pos]$ da indice del sufijo, $rank[i]$ da
posicion del sufijo $i$ en el orden.

\newpage
\section{Busqueda de subcadenas}
\textbf{Link.} \url{https://codeforces.com/group/AmklRC8lhZ/contest/695758/problem/B}

\subsection{Enunciado interpretado}
Tienes un texto fijo $S$ y consultas con patrones $P$. Para cada patron debes
decir si aparece, cuantas veces aparece o donde aparece.

\subsection{Como interpretarlo en examen}
Texto fijo + muchos patrones significa: preprocesa el texto. Con suffix array,
un patron aparece exactamente en los sufijos que empiezan con ese patron. Esos
sufijos forman un intervalo continuo en el suffix array.

\subsection{Teoria necesaria}
El orden lexicografico agrupa strings con el mismo prefijo. Si varios sufijos
empiezan con $P$, todos quedan juntos. Por eso el problema se transforma en
encontrar el primer y ultimo sufijo cuyo prefijo de longitud $|P|$ coincide con
$P$.

\subsection{Algoritmo}
\begin{enumerate}[leftmargin=*]
\item Construye $\SA$ de $S$.
\item Para cada patron $P$, haz dos binary searches:
\begin{itemize}
\item lower bound: primer sufijo $\ge P$ lexicograficamente.
\item upper bound: primer sufijo que ya no tiene prefijo $P$.
\end{itemize}
\item Si el intervalo $[L,R)$ esta vacio, no aparece.
\item Si no, aparece $R-L$ veces y sus posiciones son $\SA[L],\ldots,\SA[R-1]$.
\end{enumerate}

\subsection{Pseudocodigo corto}
\begin{lstlisting}[language=C++]
cmp_suffix_pattern(pos, P):
  // compara S[pos..] con P hasta |P|
  for i = 0..P.size()-1:
    if pos+i == n: return -1
    if S[pos+i] < P[i]: return -1
    if S[pos+i] > P[i]: return 1
  return 0  // P es prefijo del sufijo

find_interval(P):
  L = lower_bound_suffix_geq(P)
  R = first_suffix_greater_than_prefix(P)
  return [L, R)
\end{lstlisting}

\subsection{Correctitud}
Si $P$ aparece en posicion $i$, entonces el sufijo $S[i..]$ empieza con $P$.
Todos los sufijos que empiezan con $P$ tienen el mismo prefijo, asi que en orden
lexicografico forman un bloque. Las busquedas binarias encuentran los limites de
ese bloque. Por tanto el intervalo contiene exactamente todas las ocurrencias.

\subsection{Complejidad y trampas}
Construccion $\bigO(n\log n)$. Query simple $\bigO(|P|\log n)$. Trampa: si
quieres listar posiciones ordenadas por indice en el texto, el intervalo del
suffix array no necesariamente esta ordenado por posicion; debes ordenar esas
posiciones o usar otra estructura.

\newpage
\section{Contando subcadenas}
\textbf{Link.} \url{https://codeforces.com/group/AmklRC8lhZ/contest/695758/problem/C}

\subsection{Enunciado interpretado}
El problema entrena conteos con subcadenas. Segun variante, puede pedir contar
ocurrencias de patrones o contar subcadenas que cumplen alguna condicion. La
idea base es que toda subcadena es prefijo de algun sufijo.

\subsection{Como interpretarlo en examen}
Pregunta primero: ``me piden contar ocurrencias de un patron'' o ``me piden
contar subcadenas distintas''. Si es patron, usa intervalo en suffix array. Si
es distintas, usa suma de LCP.

\subsection{Teoria necesaria}
Para ocurrencias de $P$, cuenta cuantos sufijos empiezan con $P$. Para
subcadenas distintas, usa que el sufijo $S[\SA[i]..]$ aporta todos sus prefijos,
salvo los que comparte con el sufijo anterior.

\subsection{Algoritmo para ocurrencias}
\begin{enumerate}[leftmargin=*]
\item Construye $\SA$.
\item Encuentra intervalo de sufijos que tienen prefijo $P$.
\item La respuesta es el tamano del intervalo.
\end{enumerate}

\subsection{Algoritmo para distintas}
\[
total=\frac{n(n+1)}2,\qquad distinct=total-\sum_i \LCP[i].
\]

\subsection{Pseudocodigo corto}
\begin{lstlisting}[language=C++]
count_occurrences(P):
  [L, R) = find_interval(P)
  return R - L

count_distinct_substrings():
  ans = n * (n + 1) / 2
  for i = 1..n-1:
    ans -= lcp[i]
  return ans
\end{lstlisting}

\subsection{Correctitud}
Para ocurrencias, cada aparicion de $P$ corresponde a un sufijo que empieza con
$P$, y viceversa. Para distintas, cada subcadena se cuenta como prefijo de un
sufijo. En orden lexicografico, los prefijos ya vistos por el sufijo actual son
exactamente los que comparte con el sufijo anterior, medidos por LCP. Restarlos
evita duplicados.

\subsection{Complejidad y trampas}
Con $\SA+\LCP$, el conteo de distintas es $\bigO(n)$. Las consultas de patrones
son $\bigO(|P|\log n)$ en version simple. Trampa: $\sum LCP$ no cuenta
ocurrencias; cuenta prefijos duplicados entre sufijos vecinos.

\newpage
\section{Subcadenas diferentes}
\textbf{Link.} \url{https://codeforces.com/group/AmklRC8lhZ/contest/695758/problem/D}

\subsection{Enunciado interpretado}
Dada una cadena $S$, calcula cuantas subcadenas distintas tiene. Dos subcadenas
son iguales si tienen la misma secuencia de caracteres, aunque aparezcan en
posiciones distintas.

\subsection{Como interpretarlo en examen}
Esta es una pregunta casi directa de suffix array + LCP. Si ves ``distintas'',
piensa en:
\[
\frac{n(n+1)}2-\sum LCP.
\]

\subsection{Teoria necesaria}
Hay $n(n+1)/2$ subcadenas contando repetidas. Al procesar sufijos en orden de
$\SA$, el sufijo actual tiene $n-\SA[i]$ prefijos. Sus primeros $\LCP[i]$
prefijos ya aparecieron antes, porque los comparte con el sufijo anterior.

\subsection{Algoritmo}
\begin{enumerate}[leftmargin=*]
\item Construye suffix array.
\item Construye LCP array.
\item Inicializa $ans=n(n+1)/2$.
\item Resta cada valor de LCP.
\end{enumerate}

\subsection{Pseudocodigo corto}
\begin{lstlisting}[language=C++]
distinct_substrings(S):
  sa = build_sa(S)
  lcp = build_lcp(S, sa)
  ans = 1LL * n * (n + 1) / 2
  for i = 1..n-1:
    ans -= lcp[i]
  return ans
\end{lstlisting}

\subsection{Correctitud}
Cada subcadena de $S$ es prefijo de algun sufijo. En orden lexicografico, cuando
llegas a un sufijo, todos sus prefijos son candidatos a subcadenas nuevas. Pero
los prefijos de longitud hasta $\LCP[i]$ ya fueron vistos en el sufijo anterior;
los prefijos mas largos son nuevos. Por tanto el sufijo aporta
$(n-\SA[i])-\LCP[i]$ subcadenas nuevas. La suma equivale al total menos
$\sum LCP$.

\subsection{Complejidad y trampas}
Construccion $\bigO(n\log n)$ con Manber-Myers, o $\bigO(n)$ con SA-IS/DC3.
Conteo final $\bigO(n)$. Trampa: si la cadena incluye sentinela en la
construccion, no debes contar subcadenas que contengan el sentinela.

\newpage
\section{Subcadena en comun mas larga}
\textbf{Link.} \url{https://codeforces.com/group/AmklRC8lhZ/contest/695758/problem/E}

\subsection{Enunciado interpretado}
Dadas dos cadenas $A$ y $B$, encuentra la mayor longitud de una subcadena que
aparece en ambas, o la subcadena misma.

\subsection{Como interpretarlo en examen}
Dos strings + subcadena comun = concatena con separadores y usa suffix array.
No uses LCS de subsecuencia; aqui es substring, debe ser contigua.

\subsection{Teoria necesaria}
Si una subcadena $X$ aparece en ambas cadenas, entonces existe un sufijo de $A$
y un sufijo de $B$ que tienen $X$ como prefijo comun. En el suffix array de
$A\#B\$$, los sufijos con prefijo comun largo quedan juntos. Por eso basta mirar
LCP entre vecinos de origen distinto.

\subsection{Algoritmo}
\begin{enumerate}[leftmargin=*]
\item Construye $T=A\#B\$$ con separadores que no aparecen en las cadenas.
\item Construye $\SA$ y $\LCP$ de $T$.
\item Define una funcion $origin(pos)$ que diga si el sufijo empieza en $A$, en
$B$ o en separador.
\item Recorre $i=1..|T|-1$.
\item Si $\SA[i-1]$ y $\SA[i]$ vienen de cadenas distintas, actualiza con
$\LCP[i]$.
\end{enumerate}

\subsection{Pseudocodigo corto}
\begin{lstlisting}[language=C++]
longest_common_substring(A, B):
  T = A + "#" + B + "$"
  sa = build_sa(T)
  lcp = build_lcp(T, sa)
  best = 0; bestPos = -1
  for i = 1..T.size()-1:
    if origin(sa[i-1]) != origin(sa[i]):
      if lcp[i] > best:
        best = lcp[i]
        bestPos = sa[i]
  return best
\end{lstlisting}

\subsection{Correctitud}
Sea $X$ una subcadena comun optima. Entonces hay al menos un sufijo de $A$ y uno
de $B$ que empiezan con $X$. Todos los sufijos con prefijo $X$ forman un bloque
en $\SA$. Dentro de ese bloque debe existir un par adyacente de origen distinto,
porque hay sufijos de ambas cadenas. El LCP de ese par es al menos $|X|$. Por
tanto el maximo LCP entre vecinos de origen distinto recupera la respuesta
optima.

\subsection{Complejidad y trampas}
$\bigO(n\log n)$ o $\bigO(n)$ segun construccion de SA. Trampa: usa separadores
distintos y menores/no presentes; si no, puedes contar substrings que cruzan de
$A$ a $B$ artificialmente.

\newpage
\section{Ordenando subcadenas}
\textbf{Link.} \url{https://codeforces.com/group/AmklRC8lhZ/contest/695758/problem/F}

\subsection{Enunciado interpretado}
Se te da un texto $S$ y una lista de subcadenas representadas por pares
$(l,r)$. Debes compararlas, ordenarlas lexicograficamente o responder cual es
menor.

\subsection{Como interpretarlo en examen}
Comparar substrings directamente cuesta su longitud. Si hay muchas
comparaciones, necesitas LCP arbitrario entre sufijos. La estructura es:
\[
\SA + rank + LCP + RMQ.
\]

\subsection{Teoria necesaria}
Para comparar $S[l_1..r_1]$ y $S[l_2..r_2]$, primero calcula cuanto prefijo
comparten:
\[
x=lcp(l_1,l_2).
\]
Si $x$ alcanza la longitud menor, una substring es prefijo de la otra y gana la
mas corta. Si no, el primer caracter distinto esta en $l_1+x$ y $l_2+x$.

\subsection{Algoritmo}
\begin{enumerate}[leftmargin=*]
\item Construye $\SA$ y $rank$.
\item Construye $\LCP$.
\item Preprocesa sparse table sobre $\LCP$.
\item Para comparar dos substrings, usa RMQ sobre LCP para obtener LCP de los
dos sufijos.
\item Decide por longitud o por el primer caracter diferente.
\end{enumerate}

\subsection{Pseudocodigo corto}
\begin{lstlisting}[language=C++]
lcp_suffix(i, j):
  if i == j: return n - i
  a = rank[i]; b = rank[j]
  if a > b: swap(a, b)
  return rmq_lcp(a+1, b)

less_substring(l1, r1, l2, r2):
  len1 = r1-l1+1
  len2 = r2-l2+1
  x = min(lcp_suffix(l1,l2), min(len1,len2))
  if x == min(len1,len2):
    return len1 < len2
  return S[l1+x] < S[l2+x]
\end{lstlisting}

\subsection{Correctitud}
El LCP de los sufijos desde $l_1$ y $l_2$ da la longitud del prefijo comun de
las dos substrings, salvo que una termine antes. Si el prefijo comun cubre la
substring mas corta, entonces la mas corta es menor si las longitudes difieren.
Si no cubre, el primer caracter distinto decide exactamente el orden
lexicografico.

\subsection{Complejidad y trampas}
Preprocesamiento $\bigO(n\log n)$ por SA y RMQ. Cada comparacion $\bigO(1)$.
Ordenar $q$ substrings cuesta $\bigO(q\log q)$. Trampa: el LCP de sufijos puede
ser mas largo que la substring; siempre recortalo con las longitudes.

\newpage
\section{Bordes}
\textbf{Link.} \url{https://codeforces.com/group/AmklRC8lhZ/contest/695758/problem/G}

\subsection{Enunciado interpretado}
Dada una cadena, encuentra sus bordes o razona con prefijos que tambien son
sufijos. Un borde propio no puede ser la cadena completa.

\subsection{Como interpretarlo en examen}
Prefijo = sufijo es KMP/Z-function. Si el problema pide todos los bordes,
piensa en la cadena de prefix function:
\[
k=\pi[n-1],\quad k=\pi[k-1],\ldots
\]

\subsection{Teoria necesaria}
La prefix function $\pi[i]$ es la longitud del mayor borde de $S[0..i]$. Si
$\pi[n-1]=k$, entonces $S[0..k-1]$ es el mayor borde de $S$. Los bordes mas
pequenos son bordes de ese borde.

\subsection{Algoritmo}
\begin{enumerate}[leftmargin=*]
\item Calcula $\pi$ de KMP.
\item Empieza con $k=\pi[n-1]$.
\item Mientras $k>0$, registra $k$ y actualiza $k=\pi[k-1]$.
\item Invierte la lista si se pide de menor a mayor.
\end{enumerate}

\subsection{Pseudocodigo corto}
\begin{lstlisting}[language=C++]
prefix_function(S):
  pi[0] = 0
  for i = 1..n-1:
    j = pi[i-1]
    while j > 0 && S[i] != S[j]:
      j = pi[j-1]
    if S[i] == S[j]: j++
    pi[i] = j

borders(S):
  pi = prefix_function(S)
  k = pi[n-1]
  while k > 0:
    answer.push_back(k)
    k = pi[k-1]
\end{lstlisting}

\subsection{Correctitud}
Por definicion, $\pi[n-1]$ es el mayor borde de toda la cadena. Cualquier borde
menor de $S$ tambien debe ser borde del mayor borde: si coincide como prefijo y
sufijo de $S$, queda alineado dentro de ese prefijo/sufijo mayor. Por eso seguir
$\pi[k-1]$ enumera todos los bordes.

\subsection{Complejidad y trampas}
$\bigO(n)$. Trampa: revisa si el problema considera la cadena completa como
borde. Muchos problemas solo piden bordes propios.

\newpage
\section{Refran}
\textbf{Link.} \url{https://codeforces.com/group/AmklRC8lhZ/contest/695758/problem/H}

\subsection{Enunciado interpretado}
El problema trata de subcadenas repetidas. Una forma comun es buscar una
subcadena que aparece varias veces y maximiza una funcion como
$longitud \times frecuencia$, o encontrar la repeticion mas significativa.

\subsection{Como interpretarlo en examen}
Subcadena repetida = varios sufijos con prefijo comun. Varios sufijos con
prefijo comun = intervalo en suffix array. Longitud comun de un intervalo =
minimo LCP dentro del intervalo.

\subsection{Teoria necesaria}
Si tienes $t$ sufijos consecutivos en $\SA$, la longitud maxima que comparten
todos es:
\[
\min(\LCP[L+1],\ldots,\LCP[R]).
\]
Si esa longitud es $h$, entonces existe una subcadena de longitud $h$ que
aparece al menos $t=R-L+1$ veces.

\subsection{Algoritmo}
Interpreta el arreglo LCP como histograma:
\begin{itemize}[leftmargin=*]
\item una barra de altura $h$ representa que dos sufijos vecinos comparten $h$;
\item un intervalo donde todas las barras son al menos $h$ representa muchos
sufijos compartiendo un prefijo de longitud $h$;
\item si el intervalo tiene $w$ barras, hay $w+1$ sufijos.
\end{itemize}
Con una pila monotona encuentras para cada LCP el intervalo maximal donde es el
minimo. Evalua la funcion del problema.

\subsection{Pseudocodigo corto}
\begin{lstlisting}[language=C++]
best_repeated():
  sa = build_sa(S)
  lcp = build_lcp(S, sa)
  stack = empty
  best = 0
  for i = 1..n: // usa lcp[n]=0 como centinela
    start = i
    while stack not empty && stack.top.height > lcp[i]:
      (h, pos) = stack.pop()
      width = i - pos
      freq = width + 1
      best = max(best, h * freq)
      start = pos
    push (lcp[i], start)
\end{lstlisting}

\subsection{Correctitud}
Cada subcadena repetida corresponde a un conjunto de sufijos consecutivos que la
tienen como prefijo. La frecuencia es el numero de sufijos del bloque y la
longitud maxima comun es el minimo LCP del bloque. La pila monotona enumera para
cada posicion de LCP el mayor intervalo donde esa altura es el minimo. Por tanto
considera todos los bloques candidatos relevantes.

\subsection{Complejidad y trampas}
$\SA+\LCP$ domina; la pila es $\bigO(n)$. Trampa: si el problema pide devolver
la subcadena, no basta el valor; debes guardar una posicion de sufijo dentro del
intervalo ganador.

\newpage
\section{Substring periodica}
\textbf{Link.} \url{https://codeforces.com/group/AmklRC8lhZ/contest/695758/problem/I}

\subsection{Enunciado interpretado}
Se trabaja con substrings y periodos. Debes decidir si una substring puede
describirse como repeticiones de un bloque mas corto, o contar/optimizar
substrings periodicas.

\subsection{Como interpretarlo en examen}
Periodicidad significa igualdad entre una substring y ella misma desplazada.
Para muchas queries de igualdad, usa LCP arbitrario:
\[
lcp(l,l+p)\ge len-p.
\]

\subsection{Teoria necesaria}
Una cadena $X$ de longitud $m$ tiene periodo $p$ si:
\[
X[i]=X[i+p]\quad\text{para }0\le i<m-p.
\]
Eso equivale a decir que el prefijo de longitud $m-p$ coincide con el bloque
desplazado $p$ posiciones.

\subsection{Algoritmo}
Para una query substring $S[l..r]$ de longitud $m$ y candidato $p$:
\begin{enumerate}[leftmargin=*]
\item Calcula $x=lcp(l,l+p)$ con $\SA+\LCP+\RMQ$.
\item Si $x\ge m-p$, entonces $p$ es periodo.
\item Si ademas el problema pide repeticion exacta de bloques, revisa $m\bmod p=0$.
\end{enumerate}
Para buscar periodos, se prueban divisores o candidatos derivados de KMP/Z,
segun la forma exacta del problema.

\subsection{Pseudocodigo corto}
\begin{lstlisting}[language=C++]
is_period(l, r, p):
  m = r - l + 1
  if p <= 0 || p >= m: return false
  return lcp_suffix(l, l+p) >= m-p

is_exact_repetition(l, r, p):
  m = r - l + 1
  return m % p == 0 && is_period(l, r, p)
\end{lstlisting}

\subsection{Correctitud}
La definicion de periodo exige que cada caracter coincida con el que esta $p$
posiciones despues mientras ambos existan. Eso es exactamente comparar
$S[l..r-p]$ con $S[l+p..r]$. Si el LCP entre los sufijos que empiezan en $l$ y
$l+p$ alcanza $m-p$, esas dos substrings son iguales, por lo tanto $p$ es
periodo. Si no alcanza, hay una posicion que viola la definicion.

\subsection{Complejidad y trampas}
Preprocesamiento $\bigO(n\log n)$. Cada prueba de periodo $\bigO(1)$. Trampa:
``tener periodo $p$'' no siempre significa que la longitud sea multiplo de $p$;
para repeticion exacta si debes pedir divisibilidad.

\newpage
\section{Imperial Roads}
\textbf{Link.} \url{https://codeforces.com/gym/101889/problem/I}

\subsection{Enunciado interpretado}
Se tiene un grafo ponderado y se consideran consultas sobre carreteras. Para una
carretera dada, se quiere saber el costo minimo de un arbol generador que la
incluya obligatoriamente.

\subsection{Como interpretarlo en examen}
Este es un problema de MST + LCA/static tree. Primero construyes un MST. Luego
cada consulta pregunta: si fuerzo una arista $(u,v,w)$, que arista saco del
ciclo que se forma? La respuesta es sacar la arista mas pesada del camino entre
$u$ y $v$ en el MST.

\subsection{Teoria necesaria}
Propiedad del ciclo del MST: si agregas una arista a un arbol, aparece un unico
ciclo. Para volver a tener un arbol, debes eliminar una arista de ese ciclo. Si
la arista nueva es obligatoria, la mejor opcion es eliminar la mas pesada del
camino original.

Para responder maximo en camino rapido, usa binary lifting augmentado:
\[
up[v][j]=\text{ancestro }2^j,
\]
\[
mx[v][j]=\text{maximo peso en ese salto}.
\]

\subsection{Algoritmo}
\begin{enumerate}[leftmargin=*]
\item Construye MST con Kruskal y guarda su peso total $W$.
\item Enraiza el MST y preprocesa $up$ y $mx$.
\item Para consulta $(u,v,w)$, calcula $M=maxEdge(u,v)$.
\item Responde:
\[
W+w-M.
\]
\end{enumerate}

\subsection{Pseudocodigo corto}
\begin{lstlisting}[language=C++]
maxEdge(u, v):
  ans = 0
  if depth[u] < depth[v]: swap(u, v)
  lift u to depth[v], updating ans with mx[u][j]
  if u == v: return ans
  for j = LOG-1 downto 0:
    if up[u][j] != up[v][j]:
      ans = max(ans, mx[u][j], mx[v][j])
      u = up[u][j]; v = up[v][j]
  ans = max(ans, mx[u][0], mx[v][0])
  return ans

answer(edge u,v,w):
  return mstWeight + w - maxEdge(u,v)
\end{lstlisting}

\subsection{Correctitud}
Agregar la arista consultada al MST crea un ciclo formado por esa arista y el
camino $u-v$ dentro del MST. Todo arbol generador que incluya la arista
consultada debe eliminar alguna arista de ese ciclo. Para minimizar el peso
total, se elimina la arista mas pesada del camino original. La tabla $mx$ calcula
exactamente esa arista maxima al descomponer el camino en saltos de potencias de
dos.

\subsection{Complejidad y trampas}
Kruskal $\bigO(m\log m)$. LCA $\bigO(n\log n)$ memoria y preproceso. Cada query
$\bigO(\log n)$. Trampa: si la arista ya pertenece al MST, la formula sigue
dando $W$ porque el maximo del camino puede ser la misma arista con peso $w$.

\newpage
\section{Fools and Roads}
\textbf{Link.} \url{https://codeforces.com/problemset/problem/191/C}

\subsection{Enunciado interpretado}
Se da un arbol con $n$ nodos y luego muchas rutas entre pares de nodos. Para
cada arista del arbol, debes contar cuantas de esas rutas pasan por ella.

\subsection{Como interpretarlo en examen}
Muchas rutas sobre arbol + contar aristas = diferencia en arbol. No actualices
cada camino arista por arista. Marca extremos y LCA, luego acumula con DFS.

\subsection{Teoria necesaria}
En un arreglo, para sumar $1$ en un intervalo $[l,r]$, haces:
\[
diff[l]++,\quad diff[r+1]--.
\]
En arbol, una ruta $u-v$ se parte en $u\to a$ y $v\to a$, donde
$a=\lca(u,v)$. La version arborea es:
\[
cnt[u]++,\quad cnt[v]++,\quad cnt[a]-=2.
\]
Luego cada subarbol acumula hacia arriba.

\subsection{Algoritmo}
\begin{enumerate}[leftmargin=*]
\item Preprocesa LCA.
\item Para cada ruta $(u,v)$:
\[
cnt[u]++,\quad cnt[v]++,\quad cnt[\lca(u,v)]-=2.
\]
\item Haz DFS postorden desde la raiz.
\item Cuando vuelves de un hijo $x$, su acumulado final es la cantidad de rutas
que pasan por la arista $(parent[x],x)$.
\end{enumerate}

\subsection{Pseudocodigo corto}
\begin{lstlisting}[language=C++]
for each path (u, v):
  a = lca(u, v)
  cnt[u]++
  cnt[v]++
  cnt[a] -= 2

dfs(u, p):
  sum = cnt[u]
  for v in adj[u]:
    if v == p: continue
    child = dfs(v, u)
    answer_edge(u,v) = child
    sum += child
  return sum
\end{lstlisting}

\subsection{Correctitud}
Para una ruta $u-v$, los $+1$ en $u$ y $v$ suben por el arbol durante la
acumulacion. Al llegar al LCA, restar $2$ cancela esas dos contribuciones antes
de que pasen por encima del LCA. Por tanto solo quedan marcadas las aristas de
los caminos $u\to a$ y $v\to a$, que juntos forman exactamente la ruta $u-v$.
Como las contribuciones se suman linealmente, el metodo funciona para todas las
rutas a la vez.

\subsection{Complejidad y trampas}
LCA $\bigO(n\log n)$. Cada ruta $\bigO(\log n)$. DFS $\bigO(n)$. Trampa: el
valor de una arista se guarda normalmente en el hijo, no en el padre.

\newpage
\section{Frequency of String}
\textbf{Link.} \url{https://codeforces.com/problemset/problem/963/D}

\subsection{Enunciado interpretado}
Hay un texto fijo y varias consultas. Cada consulta da un patron $P$ y un numero
$k$. Se pide la menor longitud de una substring del texto que contenga al menos
$k$ ocurrencias de $P$. Si no hay suficientes ocurrencias, la respuesta es
imposible o $-1$ segun formato.

\subsection{Como interpretarlo en examen}
Primero reduce el problema: si conoces las posiciones donde aparece $P$, la
respuesta es una ventana sobre esas posiciones. La parte dificil es obtener
ocurrencias de muchos patrones. Para muchos patrones, Aho-Corasick es natural;
con suffix array tambien puedes localizar ocurrencias como intervalo.

\subsection{Teoria necesaria}
Si $P$ aparece en posiciones:
\[
p_1<p_2<\cdots<p_t,
\]
entonces una substring que contiene las ocurrencias $p_i,\ldots,p_{i+k-1}$
debe empezar no despues de $p_i$ y terminar al menos en
$p_{i+k-1}+|P|-1$. La menor ventana para ese grupo tiene longitud:
\[
p_{i+k-1}-p_i+|P|.
\]
Por tanto basta revisar ocurrencias consecutivas.

\subsection{Algoritmo con Aho-Corasick}
\begin{enumerate}[leftmargin=*]
\item Inserta todos los patrones en el trie.
\item Construye failure links.
\item Recorre el texto una vez.
\item Cada vez que un patron termina en la posicion actual, guarda su posicion
de inicio.
\item Para cada consulta $(k,P)$, toma la lista de posiciones de $P$ y calcula
el minimo de $pos[i+k-1]-pos[i]+|P|$.
\end{enumerate}

\subsection{Pseudocodigo corto}
\begin{lstlisting}[language=C++]
solve_query(pattern_id, k):
  v = occurrences[pattern_id]
  if v.size() < k: return -1
  ans = INF
  len = pattern_length[pattern_id]
  for i = 0; i+k-1 < v.size(); i++:
    ans = min(ans, v[i+k-1] - v[i] + len)
  return ans
\end{lstlisting}

\subsection{Correctitud}
Toda substring que contiene $k$ ocurrencias tiene una primera ocurrencia y una
k-esima ocurrencia dentro de ella. Si esas ocurrencias son $p_i$ y
$p_{i+k-1}$, la substring debe cubrir desde $p_i$ hasta el final de la ocurrencia
en $p_{i+k-1}$. La ventana minima para ese conjunto tiene longitud
$p_{i+k-1}-p_i+|P|$. Revisar todos los grupos consecutivos de $k$ ocurrencias
cubre todas las respuestas posibles.

\subsection{Complejidad y trampas}
Construir Aho-Corasick cuesta $\bigO(L|\Sigma|)$ o $\bigO(L)$ segun
implementacion, donde $L$ es la suma de longitudes de patrones. Recorrer texto
$\bigO(n+occ)$. Cada query cuesta proporcional al numero de ocurrencias del
patron si se hace directo. Trampa: ocurrencias pueden solaparse; no debes
descartarlas.

\newpage
\section{Ruta de estudio express para estos problemas}

\subsection{Orden recomendado}
\begin{enumerate}[leftmargin=*]
\item Static RMQ: aprende sparse table hasta poder probarla sin mirar.
\item Suffix array + LCP: entiende suffix, rank, LCP y Manber-Myers.
\item Busqueda de subcadenas: patron como intervalo en SA.
\item Subcadenas distintas y LCS: formulas clasicas con LCP.
\item Ordenando subcadenas: LCP arbitrario con RMQ.
\item Bordes: prefix function de KMP.
\item Refran/periodicidad/Frequency: problemas avanzados que combinan ideas.
\item Imperial/Fools: LCA aplicado a caminos y diferencias en arbol.
\end{enumerate}

\subsection{Plantilla universal de respuesta}
Cuando te den un problema de este bloque, responde asi:
\begin{enumerate}[leftmargin=*]
\item Identifico el objeto fijo: arreglo, arbol o string.
\item Como hay muchas consultas, preproceso.
\item Defino la estructura y que guarda.
\item Transformo la query a una operacion de esa estructura.
\item Pruebo correctitud con el invariante clave.
\item Cierro complejidad de preproceso, query y memoria.
\end{enumerate}

\subsection{Senales para no equivocarte}
\begin{itemize}[leftmargin=*]
\item ``No updates'' + rango = sparse table.
\item ``Camino en arbol'' = LCA, a veces con max/min por salto.
\item ``Muchas rutas y contar aristas'' = diferencia en arbol.
\item ``Patron en texto fijo'' = intervalo en suffix array.
\item ``Distinct substrings'' = $n(n+1)/2-\sum LCP$.
\item ``Comparar substrings'' = LCP arbitrario + RMQ.
\item ``Prefijo y sufijo'' = KMP/Z.
\item ``Muchas ocurrencias de muchos patrones'' = Aho-Corasick o SA.
\end{itemize}

\end{document}
